\section{Native endpoint tensors on ordinary tape addresses}
\label{app:native-endpoint}
Let $\Omega=\prod_{i=1}^{a}[n_i]$, $M=\prod_i n_i$,
$L=\lceil\log_2(2M)\rceil$, and $V=MD$, where $D\ge1$ is the
number of payload bits per record. The radices are even and satisfy
$L^3\le n_i\le R$, $\log R=\operatorname{polylog}L$.
Write $\tau=t$ for the supplied interchange exponent and
$\gamma=1/(2-\tau)$ for the aggregate rectangle exponent.
All coordinates below are address indices, rather than fields appended
to every record. The tape volume counts their complete Cartesian array.
\subsection{Weighted two-bank tile mask}
\label{sec:weighted-native-mask}

For an integer pair $(x,w)$, define the four-update word $P(a,b,s)$ by
the chronological operations
\[
\begin{aligned}
 x&\mathrel{+}= [w\equiv b\pmod2], &
 w&\mathrel{+}=s[x\equiv a\pmod2],\\
 x&\mathrel{-}= [w\equiv b\pmod2], &
 w&\mathrel{-}=s[x\equiv a\pmod2].
\end{aligned}
\]
The parity tests are evaluated at that step.  Set
$p=P(1,1,+1)$ and $q=P(1,0,-1)$.  A direct check on the four parity
tiles gives $p=(B\ C\ D)$ and $q=(A\ B\ D)$, with
$A=(0,0),B=(1,0),C=(0,1),D=(1,1)$.  Thus the chronological word
$p^{-1}$ followed by $q$ is
\[
 (x,w)\longmapsto
 \bigl(2\lfloor x/2\rfloor+(1-(x\bmod2)),w\bigr).
\]
For a frozen control bit $g$, mask every one of these eight elementary
updates by $g$.  The resulting word is identity for $g=0$ and the displayed
bit flip for $g=1$.  In particular it restores the whole integer $w$,
not merely its parity.  Each of the eight increments is an ordinary integer update.

Now take $u$ digit pairs with even radices $n_i,m_i$, all at least $L^3$,
and ordinary numeric orders
\[
 X=\sum_{i=0}^{u-1}x_iA_i,\quad A_i=\prod_{j<i}n_j;
 \qquad
 W=\sum_{i=0}^{u-1}w_iB_i,\quad B_i=\prod_{j<i}m_j.
\]
Each $g_i$ is computed solely from frozen address controls disjoint from
$X,W$ and the dock below; it cannot depend on a bank while that bank is
being updated.  In one $X$-update of the eight-step word, let
$\delta_i\in\{-1,0,1\}$ be its local signed increment, determined only
by $w_i\bmod2$, $g_i$, and the public word position.  Execute the bank
update by the single weighted integer increment
\[
                 \Delta_X=\sum_i A_i\delta_i.
\]
For a $W$-update, define $\epsilon_i$ from $x_i\bmod2$ and set
\[
                 \Delta_W=\sum_i B_i\epsilon_i.
\]
The first offset is independent of $X$ on each fixed $(W,\text{controls})$
fiber; the second is independent of $W$ on each fixed
$(X,\text{controls})$ fiber.  Hence each update is a cyclic line rotation
of its target numeric bank, not a gather of its component digits.

To implement it in the fixed-tape model, reserve a separate binary dock
$Z$ of width $c_0=\max(a,b)$, where
$a=\lceil\log_2\prod_i n_i\rceil$ and
$b=\lceil\log_2\prod_i m_i\rceil$.  Pad each target range once to
$2^a$ or $2^b$ (factor below two), retaining its true validity predicate.
For a selected native prefix, its numerical range is a genuine Cartesian
factor, not low binary bits of the flattened whole bank. If it is the
high factor $[N_u][N_{\rm rest}]$, append entire opposite-factor rows;
if it is the low factor $[N_{\rm rest}][N_u]$, append the invalid values
in every low-factor block. In both cases the pad is below twice the
original volume and costs one sequential scan with amortized counters.
All other coordinates are opaque spectators. Only this selected factor
is exchanged through the actual gap. For an $X$-update, interchange its $a$-bit field with an $a$-bit slice of
$Z$ using the supplied equal-chunk theorem.  On the valid prefix
$0\le X<\prod n_i$, rotate modulo the true range by $\Delta_X$; copy the
invalid tail pointwise.  If $W$ is invalid, use offset zero on the valid
$X$ prefix.  Interchange back.  The $W$-update is symmetric, forcing zero
offset when $X$ is invalid.  Every exchange restores arbitrary dirty dock
contents.  Consequently every privately padded state is fixed pointwise
by each elementary update and may be cropped only after the completed
word.  No unequal-side transpose across an arbitrary gap is presumed.
Validity follows the \emph{logical target}: after docking $X$, test
its value in the dock, not the old physical $X$ position, which now
contains arbitrary dock data. The other bank remains in place.
The internal numerical pads of each SWAP call have already been
cropped when that call returns, before this valid-prefix rotation.

Put $N_X=\prod n_i$, $N_W=\prod m_i$. Suppressing Cartesian
spectators, the original native domain is
$\Omega_0=[N_X]\times[N_W]$ and the privately padded domain is
$\widetilde\Omega=[2^a]\times[2^b]$.
Every complete docked update fixes
$\mathcal I=\widetilde\Omega\setminus\Omega_0$ pointwise.
For the concrete schedule below keep these private pads through
the repair, then crop them after the exact mask is complete.

On $\Omega_0$, define the good set $G_0$ by quotient margins
\[
 10\le\lfloor x_i/2\rfloor\le n_i/2-11,\qquad
 10\le\lfloor w_i/2\rfloor\le m_i/2-11\quad\text{for all }i.
\]
Each local integer moves by at most four during the eight updates, so on
this set no digit overflows and the weighted updates equal the independent
tile words.  The ideal completed map toggles only $x_i$'s endpoint bit;
it preserves every quotient label and fixes $W$.  It therefore preserves
$G_0$ and $\mathcal B=\Omega_0\setminus G_0$ separately.
Extend the ideal map as identity on $\mathcal I$. The packed word
$P$ agrees with that extension on $G_0\cup\mathcal I$, and both
are permutations of $\widetilde\Omega$. Hence
\[
P(\mathcal B)=\widetilde\Omega\setminus P(G_0\cup\mathcal I)
=\widetilde\Omega\setminus T(G_0\cup\mathcal I)=T(\mathcal B)=\mathcal B.
\]
The union bound gives density on the \emph{original native domain}
\[
 \delta=\frac{|\mathcal B|}{|\Omega_0|}
 \le40\sum_i(n_i^{-1}+m_i^{-1})=O(u/L^3).
\]
The repair marker is explicitly
\[
 [X<N_X]\wedge[W<N_W]\wedge[\text{a quotient margin fails}].
\]
It is false on $\mathcal I$, even if private padding occupies a
constant fraction of $\widetilde\Omega$. Do not decode an invalid
numeric bank as a wrapped native digit tuple.
Extract only marked records at their current padded output addresses $y$. Recover
their sources $x=P^{-1}(y)$ by reversing the eight packed updates, then
compute the ideal endpoint destination $T(x)$. These are public address
calculations with frozen controls; no old key is assumed to accompany
the payload. Keys are destination ordinals in the same current padded
layout. Stably sort these destination keys and refill the same bad
slots.  This corrects every exceptional address and discards no payload.
The repair costs $O(\delta VL+M_{\rm rec}L^C)$ for a fixed $C$, where
$M_{\rm rec}$ is the native record count before these private bank
pads. The padded record count and bit volume are each below four
times their native values, so full marking/refill scans cost $O(V)$.
For clarity, with $K\le\delta M_{\rm rec}$ bad records, use stable
binary radix sorting on the $b=O(L)$ key bits. In each pass copy one
fixed-width key/payload item to a staging tape, remember the selected
bit, rewind, and copy to the appropriate zero/one bucket. Concatenate
the buckets and rewind for the next pass. A fixed number of tapes
and $O(K(R_{\rm op}+b))$ steps suffice per pass, including all item
and file returns. Counter updates are amortized over the scanned
length, and resetting a counter costs at most the following item
or block length. The bound is therefore
$O(KL(R_{\rm op}+L))$, plus explicitly charged polynomial address
work and full-array marking/refill scans. At refill, consume the next
sorted item exactly at a marked bad slot;
$P(\mathcal B)=T(\mathcal B)=\mathcal B$ proves the
order without a full-key comparison at every good record. The
validity-and-margin predicate above marks $\mathcal B$ directly, so no extra mask lane is added
to the retained field codec. The exponential opaque width
below absorbs the key-only and metadata terms. The tile word uses
a constant number of dock interchanges and rotations, hence costs
$O(V(u\log R)^\tau+V)$.  In the application below, the bundled opaque
record width is exponential in $L$, so it absorbs $M_{\rm rec}L^C$; also
$\delta L=O(u/L^2)=O(1)$.  The resulting mask and its inverse cost
\[
                       O(V(u\log R)^\tau+V).
\]
After refill the padded map is exactly $T$ on $\Omega_0$ and identity
on $\mathcal I$, with the dock restored. Crop only the private $X,W$
numeric tails now, returning the original native ranges before any
enclosing phase child. A caller's Fourier dummy $p_i$ in a native
range $n_i=p_i+1$ is \emph{inside} that range and is not such a tail;
it remains part of the native kernel. Role-row, suffix and coefficient
bit pads also retain their separate outer crop boundaries below.

\subsection{Disjoint banks and available opaque lanes}
Write $x_i=2q_i+b_i$. The desired local matrix is $H$ on
$\{2q_i,2q_i+1\}$ for every $q_i$; hence it is $H^{\otimes a}$
on the endpoint bits at fixed quotient labels. This is an identity
on fibers of the existing native array, not an uncharged regrouping.

Put $\ell_R=\lceil\log_2R\rceil$ and
\[
 h=\min\{\lfloor a/16\rfloor,\lfloor L/(40\ell_R)\rfloor\}.
\]
For each selected block of at most $2h$ axes, one aggregate rectangle
move places its native interval first. Take the selected endpoint
planes as $X,Y$, and three disjoint $h$-axis banks from the other
axes as $Q,Z_0,W$. The remaining axes form a spectator suffix.
The $j$th digit is $2X_j+Y_j$; the two banks need not be physically
interleaved. The banks $Q,Z_0$ are dirty selectors, and $W$ is the
separate tile auxiliary. In a shear updating any of these banks,
its controls exclude the target, $W$ and the binary dock. In particular
the source-value toggle of $Q$ reads $X,Y$, the toggle of $Z_0$
reads $X,Y,Q$, and a target transposition reads the other endpoint
plane and $Z_0$. Each completed echo restores both dirty selectors.

Pad the whole spectator suffix once to its next binary range, with
factor below two. Its index width $s$ satisfies $s\ge7L/8-O(1)$,
because $5h\log_2R\le L/8$. For the fixed final body of
Lemma~\ref{lem:body}, including its full finite phase cover, reserve
\[
 r=\left\lceil\log_2 B^{\lceil\log_N(h+2)\rceil}\right\rceil
       =O_N(\log(h+2))
\]
high suffix bits for independent numeric rows. Pad that row range
once to a multiple of $B^{\lceil\log_N(h+2)\rceil}$; the factor is
below two. Reserve the next
$c_0=\max_G\lceil\log_2|G|\rceil\le h\log_2R+1$
bits as a dock for $G\in\{X,Y,Q,Z_0,W\}$. The low suffix and
original payload together have opaque width
\[
 R_{\rm op}=2^{s-r-c_0}D\ge2^{4L/5}D
\]
above an effective fixed cutoff. Move only the $r$ high row bits to
the front, at $O(Vr)$ cost, once per completed block. After the row
pad the numeric range need not be binary; its $r+1$-bit counter is
metadata, not a further borrowed coordinate. Crop the numeric row
pad before undoing these digit moves, then crop the whole suffix
only after the completed block has fixed it.

Here and below an aggregate rectangle is padded as a whole, exchanged
by Lemma~\ref{lem:movement}, and cropped after its full endpoint.
The volume factor is below four. No individual native axis is padded
to a power of two. For a short selected block use the same available
banks at its smaller width. Since $h\le a/16$, its required auxiliary
axes are disjoint from the selected ones.

\subsection{Address generation and exact sparse repair}
The instruction descriptors consist of fixed body-slot numbers, active
rank lists, disjoint matchings, group placements, and a list of
$O_N(\log h)$ inherited mask predicates. Decode the five bank indices
by schoolbook division on $O(L)$-bit integers. Store rank lists in
$O(L)$ records of $O(L)$ bits; a lookup is a sequential scan, not
unit-cost indexing. A fixed-slot row operation pairs ranks
$(i-1)f+j$ and $(k-1)f+j$ for $1\le j\le f$. Interval placements
are obtained by tracing their cycles and using two cycle reflections.
The loops have bounds $O_N(L)$ or $O_N(L^2)$ and use binary arithmetic
of polynomial cost. Thus one fixed polynomial $C_NL^c$ bounds the
generation of any current key, selector, inverse packed update and
source scalar. No expanding Boolean expression is substituted for
the list of inherited predicates.

The width $R_{\rm op}$ dominates this polynomial beyond an effective
fixed cutoff, even when $D=1$. With $K$ exceptional opaque records
and $b=O(L)$ key bits, a staged stable radix sort costs
$O(Kb(R_{\rm op}+b))$, plus $O(V)$ full marking/refill scans and
$M_{\rm rec}C_NL^c$ metadata. The tile bad fraction is $O(u/L^3)$,
so its payload repair is $O(Vu/L^2)=O(V)$. The active-prefix bad
fraction below is at most $(L+2)^{-128}$ by increasing its fixed
prefix constant. Its full-payload repair is also $O(V)$. Counter
resets, bucket rewinds, erasure and head returns are included in
these scans. Repair keys are computed from output addresses by
the inverse packed map, never by an unrecorded input pointer.

\subsection{Restored selectors and active-rank permutations}An $A$-digit has states $0,1,2$ active and state $3$ idle. Represent
$d_i=2X_i+Y_i$. For any active mask $S$, $A^{\otimes h}$ restricts to
$C^{\otimes |S|}$ on the active coordinates and identity on idle
coordinates. All $4^h$ addresses remain in the stream; no sector is
gathered out.

The mask-routing identities use dirty address fields. Let $T_Z$ apply
one fixed transposition of active symbols where dirty selector bit
$Z_i=1$. It preserves the active mask $S$, but arbitrary dirty $Z$
can make it change paired sources as well as targets. Source-value
predicates are recomputed only after the completed $T_{J\wedge Q}$
word below, which changes matched targets only. For a mask-only selector
$J(S)$ the chronological word
\begin{equation}
 Z\mathrel{\oplus}=J(S),\quad T_Z,\quad
 Z\mathrel{\oplus}=J(S),\quad T_Z                         \label{eq:router-11}
\end{equation}
restores arbitrary initial $Z$ and equals $T_J$. With a second dirty
field $Q$, use $J(S)\wedge Q$ in \eqref{eq:router-11}. If $F_a(d)_i$ says the paired
source of target $i$ currently has trit $a$, then the completed
$T_{J\wedge Q}$ changes targets only and preserves all paired sources.
Consequently
\begin{equation}
 Q\mathrel{\oplus}=F_a(d),\quad T_{J\wedge Q},\quad
 Q\mathrel{\oplus}=F_a(d),\quad T_{J\wedge Q}              \label{eq:router-12}
\end{equation}
restores arbitrary $Q,Z$ and implements the target transposition
conditioned on source value $a$. The source/target matching must be
disjoint; this is what makes the second $F_a$ equal to the first at the
completed inner-operation boundary. A third dirty field is reserved for
the weighted tile auxiliary. The supplied weighted-mask lemma implements
each frozen XOR in \eqref{eq:router-11}--\eqref{eq:router-12} with its separate auxiliary restored;
its exact sparse boundary repair and interchange cost are not reproved
here.

The two active transpositions have explicit bit circuits
\begin{equation}
 (0\ 1):\ Y\mathrel{\oplus}=Z\wedge\neg X,
 \qquad
 (1\ 2):\ X\mathrel{\oplus}=Z\wedge Y;\quad
 Y\mathrel{\oplus}=Z\wedge X;\quad
 X\mathrel{\oplus}=Z\wedge Y.                           \label{eq:router-13}
\end{equation}
They are constant sequences of frozen-control XOR shears. Applying the
two transpositions in chronological order $(1\ 2)$ then $(0\ 1)$ adds
one to an active trit; reverse order subtracts one. Conditioning on
source trit and composing these shears realizes a matching SUM
\begin{equation}
 \text{target}_i\gets\text{target}_i+\text{source}_i\pmod3.
                                                               \label{eq:router-14}
\end{equation}
Three SUMs plus a negation implement a coordinate swap. All dirty
selectors and the tile auxiliary return to their arbitrary incoming
values after each completed operation.

For a node acting on the first $e$ active ranks, restrict controls and
endpoints to the first
\[
 u=\min\{h,\lceil C_0(e+\log(L+2))\rceil\}
\]
physical positions, where $L\ge h$ is the ambient address-width bound
in the native implementation. If fewer than $e$ active positions occur there,
define the operation to be identity; otherwise all required ranks are
present. Restrict every raw $T_Z$ transposition and both selector-bank
toggles to this prefix, not just the completed selector $J$. Insert
inherited eligibility into $J(S)$; it is not a frozen control on an
inner bit shear whose target can change $S$ temporarily. Individual
inner shears still read only the other bit plane and the selector bank.
Their completed transposition preserves $S$. The comparison with the
full intended route is made only after a \emph{complete logical}
mask-selected transposition, source-conditioned operation \eqref{eq:router-12}, SUM
or recursive basis placement on the first $e$ active ranks, with all
its dirty selectors restored. Write this
prefix implementation as $P_u$ and its intended permutation as $T$.
On a good eligible prefix all intended endpoints are inside it, so
$P_u=T$. On a bad or ineligible prefix the selector $J$ is zero:
the two raw transpositions in \eqref{eq:router-11} cancel, each inner operation is
identity, and the two source-predicate toggles in \eqref{eq:router-12} cancel too.
Thus the completed $P_u$ is identity there. These statements do not
require raw dirty $T_Z$ to preserve paired sources.

One must not repair an internal truncated $T_Z$ to its full-$h$
version: with $h=2,u=e=1$, $d=(0,0)$ and $Z=(0,1)$, the prefix is
good but full $T_Z$ for $(0\ 1)$ changes the second digit and truncated
$T_Z$ does not. This toy example embeds for any actual $u<h$ by
putting the dirty selector immediately after the prefix and
using active digits throughout the prefix.
That off-prefix action is intended to cancel in the
outer word. Nor is the prefix bad set invariant after an inner bit
shear: the first shear of $(1\ 2)$ sends idle $(X,Y)=(1,1)$ to active
$(0,1)$. Its tile repair uses a different invariant, as proved in the tile construction above.
The separate top-level digit-group placement can reach ranks beyond
$e$; it retains its full-$h$ root-level charge and is not justified
by this width-$e$ prefix argument.
In the uniform $[4]^h$
address space the prefix active count is
binomial with mean $3u/4$. For $u\ge2e$ and $u<h$, a Chernoff bound gives
\begin{equation}
 \Pr\{B_u<e\}\le e^{-u/24}.                            \label{eq:router-15}
\end{equation}
For $u=h$, eligibility $|S|\ge e$ leaves no exceptional eligible mask.
Both $T$ and the completed logical $P_u$ preserve $S$, so the
exceptional set
$E_{\rm inherited}(S)\wedge[|S|\ge e]\wedge[B_u<e]$ is invariant.
Its whole-universe fraction is bounded by the unconditional binomial
estimate; no conditional distribution given eligibility is assumed.
At that completed boundary, sorting the exceptional records by
$T(P_u^{-1}(y))$ and refilling their slots repairs it exactly;
the maps agree everywhere else. No sector is dropped. The full role recursion always
stores the complete $[4]^h$ address universe, and its row split depends
only on the independent role field, never on $S$. Thus every child
retains the same mask distribution and off-eligibility maps are identity.
Consequently, on every sector $S$, the role program applies
$C^{\otimes |S|}$ to its active digits and identity to its idle digits;
the direct sum over all masks is exactly $A^{\otimes h}$.

The prefix correction above belongs to the completed logical routing
permutation, including $T$ and $T^{-1}$ below, after all its source
conditioning and before a linear combination forms a phase. An
internal $T_{J\wedge Q}$ is not separately corrected when its surrounding
source predicate was prefix-truncated too. It never sorts the output
of the non-permutation matrix $C$.

At selected width one, let $T$ cyclically permute that eligible active
trit and be identity on ineligible masks. The mask-only transpositions
above implement $T$ and $T^{-1}$. Three temporary array copies and a
pointwise linear combination compute
$I+\eta^2(T+T^{-1})$, which is $C$ on the selected trit. Off eligibility
the two identical terms cancel, leaving the identity. The copies,
returns and cleanup add $O(V)$ at this base call. Norm two uses its
inverse via the fixed coefficient wrapper described below.

\subsection{The paired phase on the full address universe}
For $e$ paired coordinates, keep the first $2e$ active trits in
two ordered halves. Every child inherits its parent's mask predicate;
its domain test $|S|\ge2e$ is conjoined with that predicate. A
prefix failure is repaired at the completed logical routing boundary
as above, rather than permanently disabling the ideal phase.

The rank-one kernel has an explicit nine-term formula. Summing the
inverse Fourier expression first over one frequency gives
\[
 K_{I_1}=\sum_{a,b\in\F_3}\eta^{-ab}T_{(a,b)},\qquad
 (T_{(a,b)}f)(x,y)=f(x+a,y+b).
\]
Indeed $\sum_{\zeta}\eta^{\zeta(\xi-y)}$ is one for $\xi=y$ and
zero otherwise in characteristic two. The coefficients sum to one.
Use nine owned copies, the completed selected translations and a
pointwise linear combination. Off the inherited gate every
translation is identity, so the same formula is identity there.
The copies, returns and erasure cost $O(V)$.

Freeze the final body with $a_*=P/B<N^z$ for a fixed $z<\tau$.
Its charts are finite products of ternary row additions, negations
and swaps, lifted by the column identity. The restored selector
permutations above implement these exact disjoint-slot matchings.
For $e=Nf+r$, $0\le r<N$, first align the $Nf$ coordinates in
each half, keeping both $r$-coordinate remainders together. This is
a permutation of the original first $2e$ active ranks, implemented
by two reflections of its cycles. Execute the body on the $Nf$
pairs, place and process the $r$ remaining pairs by bounded direct
kernels, and undo the alignment.

The independent row stock was prepared once. Each child consumes
one factor $B$ and receives exactly $V/B$ bits, including every
mask, spectator, codec lane and retained root pad. Each of the $P$
queries, including owned queries, is charged on that volume. Owned
initialization, scalar maps and cleanup add $O_N(V)$ at the node.
Before a child, push only the other current sibling streams onto a
parking stack and copy its active stream. On return pop those
siblings from right to left into their saved intervals. The parent
boundary is marked; no descendant head crosses into a parked
ancestor. A fixed set of input/output, role, owned, padding/repair
and stack tapes suffices. Each of the fixed $P$ push/pop operations
costs $O_N(V)$ at its parent, not a multiplier on descendant time.

Consequently the time per active input bit satisfies, for one fixed
polynomial $Q$ in $\log(L+2)$,
\[
 F(e;L)\le a_*F(\lfloor e/N\rfloor;L)
       +Q(L)((e+\log(L+2))^\tau+Q(L)).
\]
At fixed bounded $e$ use the direct kernels. Unrolling the $e$-dependent
toll has ratio $a_*/N^\tau<1$. Width-independent tolls total
$O((e+1)^z\operatorname{polylog}L)$, with an extra logarithm if
$a_*=1$. Thus
\[
 F(e;L)=O((e+1)^\tau\operatorname{polylog}L).
\]
No child repeats root layout, padding or lane encoding, and each
temporary selected-factor pad is cropped before entering a child.

On four active trits, the chart $J$ in Lemma~\ref{lem:endpoint}
gives
$U_{J^T}C^{\otimes4}U_{J^T}^{-1}=K_{I_2}$.
For each mask with $b$ active trits, write $\lfloor b/4\rfloor$ in
base $N$. Its digits specify at most $N-1$ consecutive groups of
$4N^j$ active ranks per level. A group placement is two cycle
reflections; it costs $O(V(h\log R)^\tau\operatorname{polylog}L)$.
Apply the chart, one paired rank-$2N^j$ job, its inverse chart and
the inverse placement. Process the at most three remaining trits
by direct $C$ kernels. There are $O_N(\log(h+2))$ levels and the
geometric sum of paired costs is $O(V(h+1)^\tau\operatorname{polylog}L)$.
All grouping predicates depend only on the mask, so their description
does not enumerate its exponentially many values. The completed map
on every mask is $C^{\otimes b}$ on its active trits and identity on
its idle trits. The direct sum is exactly $(C\oplus[1])^{\otimes h}$.

\subsection{Completion of the binary endpoint tensor}
Appendix~\ref{app:positive-word} gives a fixed word of $384$ forward
phase letters and $392$ monomial letters over $\F_{16}$ whose matrix
is $H\otimes H$. Tensor its letters across the selected $h$ pairs,
rather than running a separate word per axis. Every monomial
permutation on four states is affine on two binary bits and uses
the weighted masks. If $M_{\pi(d),d}=m_d$, generate
$\prod_jm_{d_j}$ at the source address, apply it and then route by
$\pi$; the source-indexed scalar list is not reused at the output
address. Its generator has the polynomial cost paid above.

Encode the $D$ payload bits once into groups of four binary coefficient
planes, with at most three zero bits of codec padding per record. The
volume factor is at most four and is paid only at block entry.
The four actual coefficient planes and all borrowed rows, docks and
opaque lanes survive the complete word. The completed matrix is
binary and fixes the borrowed coordinates, so zero codec padding
returns to zero and may then be cropped. The high row pad crops
before undoing its bit moves; the suffix pad crops afterwards.
Private mask pads cropped at their earlier completed boundaries do
not propagate into the phase recurrence. All crops and returns have
their charged scans.

Since $a\le L/(3\log_2 L)$, blocks of at most $2h$ cover the axes
in $\operatorname{polylog}L$ passes: either $h\ge a/32$ or
$h\ge L/(80\ell_R)$. The finite remaining range has a direct
branch. Aggregate layouts cost $O(VL^\gamma)$ per block; phase and
monomial work costs $O(VL^\tau\operatorname{polylog}L)$. Because
$\gamma>\tau$, the total is
$O(VL^\gamma\operatorname{polylog}L)$, and in particular the bound
stated in Lemma~\ref{lem:endpoint}. The completed ordinary bit
operator is symmetric. The same procedure computes its literal bit
transpose; two applications compute its inverse because $H^3=I$.
